\section{Experiments and Results}
\label{sec:experiments_results}

The five-centre cohort provides a clinically realistic setting in which multimodal imaging evidence is abundant but report-level semantic supervision is incomplete. This imbalance is central to the evaluation of LCAD--RASA, because the proposed framework uses structured report semantics as an alignment signal without assuming that physician-authored reports are available for every case. The results are therefore organised as a single evidentiary sequence: the cohort first defines the report-supervision bottleneck; pseudo-report analyses then test whether missing reports can be converted into modality-grounded semantic targets; alignment and scarcity analyses examine whether these targets behave as useful supervision; held-out and external-baseline evaluations test predictive behaviour under a locked protocol; and ablation, perturbation, centre-shift, safety, and scalability analyses define the claim boundary of the study.

\subsection{Cohort scale and report-supervision imbalance}
\label{subsec:cohort_results}

The final analytic cohort consisted of 1,897 multimodal cervical screening cases from five centres, containing 137,591 evaluable OCT and colposcopy images (Table~\ref{tab:cohort_summary}, ``Cohort scale and case-level report supervision''; Figure~\ref{fig:centre_supervision}, ``Centre-level cohort scale and report-supervision imbalance''). Thus, image-level evidence was available at sufficient scale for multimodal representation learning. In contrast, structured clinician-authored report supervision was substantially less complete: archived real reports were available for 744 cases, whereas 1,153 cases had complete multimodal evidence but no archived report and were therefore eligible for pseudo-report completion.

The report-supervision imbalance was strongly centre-dependent. Enshi retained real reports for all 406 cases, Jingzhou retained reports for 334 of 406 cases, Xiangyang retained only 4 reports among 500 cases, and Wuhan and Shiyan had no centre-wide real-report archives. This distribution shows that the principal bottleneck was not imaging-scale data acquisition, but the uneven availability of semantic supervision. It motivates a report-aware learning strategy that can use real reports where available while extending semantic alignment to cases with complete multimodal evidence but missing report records.

\subsection{Modality-grounded pseudo-report generation}
\label{subsec:pseudo_report_results}

The pseudo-report analysis was designed to distinguish structured semantic augmentation from label-to-text expansion. The label-template baseline provided a negative control: although it satisfied the expected report schema, it had no modality support, with a mean modality-support rate of 0.000, and showed marked textual repetition, with a maximum duplicate fraction of 0.867 (Table~\ref{tab:pseudo_source_comparison}, ``Structured pseudo-report source comparison''; Figure~\ref{fig:pseudo_source_comparison}, ``Structured pseudo-report source comparison across template, rule-based, and local LLM sources''). Schema validity alone was therefore insufficient to establish useful report supervision, because a syntactically complete report could still fail to encode OCT, colposcopy, or clinical evidence.

Rule-based and local embedding-assisted LLM pseudo reports showed a different pattern. Both preserved OCT, colposcopy, and clinical section support, with a mean modality-support rate of 1.000. The rule-based baseline reached a label-consistency mean of 0.628, and the local LLM baseline reached the same mean while retaining modality-grounded sections. These results indicate that structured pseudo reports can supply section-level semantic anchors linked to the available modalities, rather than merely restating diagnostic labels in textual form.

The external-provider comparison further assessed whether pseudo-report generation depended on a single LLM source. Template, rule-based, local embedding LLM, Qwen-Plus, GLM-4.7-Flash, Gemini-3.1-Pro, and GPT-5.5 were evaluated under a common structured-report schema (Table~\ref{tab:llm_provider_comparison}, ``LLM provider comparison for structured pseudo-report generation''; Figure~\ref{fig:api_provider_summary_heatmap}, ``LLM provider comparison for structured pseudo-report generation''; Figure~\ref{fig:api_quality_heatmap}, ``Stage 1 provider quality metrics and risk-oriented quality indicators''; Figure~\ref{fig:api_reliability}, ``Provider latency, modality support, and generation reliability''). All evaluated sources produced schema-valid structured outputs in their available cohorts. In the extended 100-case cohort, GPT-5.5 generated 100 valid structured reports, with a schema-valid rate of 1.000, section completeness of 0.810, modality support of 0.810, contradiction rate of 0.050, hallucination rate of 0.000, mean latency of 11.39 seconds per case, and an estimated cost of US\$27.73 per 1,000 cases. Qwen-Plus, GLM-4.7-Flash, and Gemini-3.1-Pro were retained as pilot cohorts; their results are therefore used as feasibility evidence rather than as a definitive provider ranking.

For all external API evaluations, only de-identified structured evidence was submitted. Raw images, image paths, patient names, internal patient identifiers, and hospital identifiers were not transmitted. This design preserves the role of pseudo reports as a controlled learning component and avoids treating API-generated text as a replacement for physician-authored clinical documentation.

\subsection{Section-wise semantic alignment of multimodal representations}
\label{subsec:alignment_results}

Having established that pseudo reports retained modality-specific content, the next analysis evaluated whether report sections could function as explicit alignment targets. OCT embeddings were matched to \texttt{oct\_findings}, colposcopy embeddings to \texttt{colposcopy\_findings}, clinical embeddings to \texttt{clinical\_context}, and fused representations to \texttt{impression}. Full LCAD--RASA achieved a macro mean reciprocal rank (MRR) of 0.051 across the four sections, whereas removing section alignment reduced macro MRR to 0.021 and simple concatenation fusion remained low at 0.022 (Table~\ref{tab:theme1_alignment_ablation}, ``Direct RASA alignment ablation using modality-section retrieval''; Figure~\ref{fig:alignment_retrieval_mrr}, ``Modality-section retrieval MRR for Theme 1 cross-modal semantic alignment'').

The real-report-only variant achieved the highest retrieval-only macro MRR, at 0.057. This finding is consistent with the role of physician-authored reports as strong semantic targets when such reports are available. However, because real reports were sparse and centre-imbalanced, the retrieval advantage of the real-report-only variant does not by itself solve the supervision-coverage problem. LCAD--RASA addresses this distinction by extending report-guided alignment beyond the subset of cases with archived reports, while preserving a clear separation between real-report supervision and pseudo-report augmentation.

The staged API alignment analysis showed that section-level alignment remained measurable across pseudo-report sources. Gemini-3.1-Pro achieved the highest available pilot macro MRR of 0.366, whereas GPT-5.5 achieved a macro MRR of 0.063 on the extended cohort (Figure~\ref{fig:api_alignment_mrr}, ``Stage 2 provider alignment analysis: macro MRR and section-specific MRR''; Figure~\ref{fig:api_section_mrr}, ``Provider-level section MRR across report sections''). Because cohort size and evaluation stage differed across providers, these values should not be interpreted as a definitive provider ranking. They instead demonstrate that structured report sections provide tractable alignment targets whose behaviour can be quantified across generation sources.

\subsection{Effect of pseudo-report augmentation under report scarcity}
\label{subsec:scarcity_results}

The relevance of pseudo-report augmentation was most evident when real-report supervision was deliberately restricted. In the report-supervision scarcity curve, the largest improvement appeared in the lowest-supervision regime (Table~\ref{tab:scarcity_curve}, ``Report-supervision scarcity curve using a lightweight downstream surrogate''; Figure~\ref{fig:scarcity_curve}, ``Report-supervision scarcity curve comparing real-report-only and pseudo-report-augmented surrogates''). At 10\% real-report availability, the LCAD-augmented surrogate achieved AUROC 0.808 and F1 0.500, compared with AUROC 0.630 and F1 0.316 for the real-report-only surrogate. At full real-report availability, the augmented surrogate also remained higher, with AUROC 0.838 and F1 0.542 compared with AUROC 0.799 and F1 0.471 for the real-report-only surrogate.

This pattern suggests that pseudo-report augmentation contributes most when images and structured clinical variables are available but semantic report supervision is incomplete. The gain is therefore aligned with the cohort-level bottleneck identified above, rather than reflecting a simple increase in textual volume. The staged API downstream surrogate was intentionally conservative: GPT-5.5 achieved AUROC 0.630 and F1 0.229 at 10\% real-report availability on the extended cohort, whereas Gemini-3.1-Pro achieved AUROC 0.725 and F1 0.379 in its pilot cohort (Figure~\ref{fig:api_scarcity_auc}, ``Stage 3 downstream scarcity surrogate for selected API-generated pseudo reports''). These analyses support the utility of structured pseudo-report augmentation under report scarcity, while remaining complementary to full end-to-end retraining evidence.

\subsection{Held-out diagnostic performance}
\label{subsec:diagnostic_results}

The held-out test analysis evaluated whether the semantic-alignment pathway translated into diagnostic prediction under a locked retrospective protocol. On the held-out test set ($n=288$), Full LCAD--RASA achieved AUROC 0.832 (95\% CI 0.757--0.897) and F1 0.611 (95\% CI 0.458--0.737) at the validation-selected threshold of 0.37 (Table~\ref{tab:main_comparison}, ``Held-out test performance at validation-selected thresholds''; Figure~\ref{fig:main_auc}, ``Held-out AUROC with bootstrap confidence intervals''; Figure~\ref{fig:main_metrics}, ``Held-out multi-metric profile across model variants''; Figure~\ref{fig:main_auc_f1}, ``AUROC--F1 trade-off across main model variants''). Pseudo-augmented LCAD training achieved AUROC 0.826 (95\% CI 0.747--0.894) and F1 0.545 (95\% CI 0.395--0.674). Real-report-only training achieved AUROC 0.725, whereas simple concatenation fusion achieved AUROC 0.472.

The ordering of the internal model variants was consistent with the preceding analyses. Generic concatenation was insufficient for the multimodal setting; real-report-only training was constrained by incomplete report coverage; pseudo-augmented training improved over real-report-only learning; and the full LCAD--RASA configuration achieved the highest point estimates among the report-aware internal variants. Corrected paired bootstrap rechecking supported clear AUROC differences between Full LCAD--RASA and weaker internal comparators such as simple concatenation fusion and real-report-only training, but not between Full LCAD--RASA and pseudo-augmented LCAD. Accordingly, these results should be interpreted as favourable internal held-out performance under the locked retrospective protocol, rather than as evidence of clinical superiority.

\subsection{External same-split baseline comparison}
\label{subsec:external_baseline_results}

To address whether the proposed framework was being compared only against weak internal variants, we added an external baseline block under the same train/validation/test split, validation-selected threshold rule, held-out test set, and bootstrap confidence-interval protocol. The block included clinical-only logistic regression, clinical-only gradient boosting, OCT-only and colposcopy-only frozen-embedding MLP classifiers, a late-fusion MLP, a cross-attention multimodal transformer, and a CLIP-style contrastive multimodal baseline without report-section supervision (Table~\ref{tab:external_baselines}, ``Same-split external baseline comparison on the locked held-out test set''; Figure~\ref{fig:external_baselines_auc}, ``Same-split external baseline AUROC comparison with bootstrap confidence intervals''; Figure~\ref{fig:external_baselines_metrics}, ``Metric profile for Full LCAD--RASA and external baselines under the locked held-out protocol''). Clinical-only baselines used age, HPV, and TCT only; centre identifiers, report text, pathology-like attributes, and labels were excluded from their inputs.

This comparison showed that Full LCAD--RASA was competitive with several conventional baselines, including clinical-only logistic regression (AUROC 0.830), colposcopy-only embedding MLP (AUROC 0.827), and late-fusion MLP (AUROC 0.824). However, it did not dominate all stronger external baselines. Clinical-only HistGradientBoosting reached AUROC 0.862, the cross-attention multimodal transformer reached AUROC 0.847, and the CLIP-style contrastive multimodal baseline reached AUROC 0.888. In the corrected paired bootstrap recheck, the contrastive baseline exceeded Full LCAD--RASA in AUROC by 0.056 (95\% bootstrap interval 0.024--0.092 in favour of the comparator; two-sided $p=0.002$), whereas differences against clinical-only logistic regression, clinical-only HistGradientBoosting, late-fusion MLP, and cross-attention transformer were not statistically conclusive (Table~\ref{tab:paired_bootstrap_recheck}, ``Corrected paired bootstrap AUROC recheck''; Figure~\ref{fig:external_pairwise_recheck}, ``Corrected paired bootstrap AUROC differences between Full LCAD--RASA and external baselines'').

These results sharpen the claim boundary of the study. LCAD--RASA should not be presented as a universally superior pure discriminative classifier over all standard multimodal baselines. Its main contribution is instead report-aware multimodal learning: it converts incomplete report supervision into structured, modality-grounded semantic anchors; provides section-level alignment targets; and supports auditable pseudo-report generation under sparse report supervision. The strong contrastive result also motivates future work combining report-section supervision with contrastive multimodal pretraining.

\subsection{Contribution of report-aware alignment components}
\label{subsec:ablation_results}

Ablation analyses linked the held-out performance pattern to the report-aware alignment design. Removing section alignment reduced direct alignment quality, with macro MRR decreasing from 0.051 for Full LCAD--RASA to 0.021 (Table~\ref{tab:theme1_alignment_ablation}, ``Direct RASA alignment ablation using modality-section retrieval''). In the held-out comparison, the LCAD variant without section alignment retained a similar AUROC of 0.826, but its F1 score decreased to 0.510 compared with 0.611 for Full LCAD--RASA (Table~\ref{tab:main_comparison}, ``Held-out test performance at validation-selected thresholds''). This discrepancy indicates that section-level alignment may influence threshold-dependent error balance and operating-point behaviour even when rank-based discrimination remains similar.

Additional modality, quality-control, alignment-weight, and component ablations are reported in Supplementary Tables~S1 and S3--S5 and Figure~\ref{fig:rasa_lambda} (``Alignment-weight sweep for the RASA objective''), Figure~\ref{fig:rasa_component_boxenplot} (``Modality ablation and RASA component ablation''), and Figure~\ref{fig:qc_ablation} (``LCAD quality-control ablation''). These supplementary ablations were used as diagnostic analyses rather than as final-model superiority tests, because several variants were run under fixed quick-budget settings and some variants produced point estimates close to the full model. The QC ablation in particular did not materially change downstream AUROC/F1 under the current retrospective configuration, so QC is treated here as a procedural safeguard for pseudo-report generation rather than as a major independent driver of diagnostic performance. Collectively, the ablations support the internal coherence of the RASA design, but they should not be interpreted as causal proof that any single component independently determines the full diagnostic behaviour.

\subsection{Perturbation-based evidence for modality-grounded decoding}
\label{subsec:perturbation_results}

Perturbation analysis further examined whether generated report sections depended on their corresponding modality evidence. Removal or disruption of a modality produced the largest degradation in the matched report section (Table~\ref{tab:perturbation_matrix}, ``Upgraded perturbation sensitivity matrix''; Figure~\ref{fig:perturbation_heatmap}, ``Perturbation condition by report-section similarity matrix''; Figure~\ref{fig:perturbation_lineplot}, ``Section-level degradation profiles under modality perturbations''; Figure~\ref{fig:risk_delta}, ``Risk-score shifts under report-generation perturbations''; Figure~\ref{fig:theme1_perturbation}, ``Theme 1 upgraded perturbation sensitivity matrix''). Masking OCT produced a 0.743 drop in the OCT findings section, masking colposcopy produced a 1.000 drop in the colposcopy findings section, and masking clinical instruction produced a 0.833 drop in the clinical-context section.

Label-only inference produced the largest overall report drop of 0.466 and the largest absolute risk shift of 0.371. The degradation profile indicates that report generation was not fully recoverable from class labels alone and that section content remained sensitive to modality-specific evidence. This supports the interpretation that LCAD--RASA performs modality-grounded decoding rather than fixed-template schema filling. The evidence is perturbational rather than causal, and generated reports should therefore not be interpreted as causal clinical explanations.

\subsection{Reference-stratified performance and centre-level generalisation}
\label{subsec:loco_results}

Reference-stratified testing assessed whether held-out performance was confined to cases with archived reports. Full LCAD--RASA achieved AUROC 0.824 among cases with reference reports ($n=108$) and AUROC 0.850 among cases without reference reports ($n=180$). The comparable performance across these strata suggests that the model's held-out behaviour was not restricted to the real-report subset, although the analysis remains retrospective and conditioned on the available cohort composition.

Cross-centre evaluation provided a stricter assessment of generalisation. Under strict leave-one-centre-out retraining with the fixed quick-budget protocol, held-out AUROC was 0.702 for Enshi, 0.648 for Jingzhou, and 0.382 for Xiangyang (Supplementary Tables~S2 and S2b; Figure~\ref{fig:loco_heatmap}, ``Leave-one-centre-out evaluation: AUROC heatmap and forest-style centre summary''; Figure~\ref{fig:loco_forest}, ``Strict LOCO centre-wise behaviour under fixed quick-budget retraining''). Eval-only LOCO using a global checkpoint is reported separately and is not pooled with strict LOCO retraining. These results expose centre shift as a material limitation of the current model. The evidence therefore supports report-aware fusion in retrospective held-out testing, but does not justify an unrestricted external-validation claim across all centres.

\subsection{Robustness, safety, and computational feasibility}
\label{subsec:robustness_safety_scalability}

Robustness and feasibility analyses examined whether the pipeline remained reproducible and computationally tractable at cohort scale. The supplementary evaluation included masking validation, modality ablations, quality-control ablations, multi-seed stability, report-level safety metrics, and runtime summaries (Supplementary Tables~S7--S11; Supplementary Figure~\ref{fig:supplementary_robustness}, ``Supplementary robustness checks: masking validation and multi-seed stability''). The pipeline processed approximately 137k images, and the publishable model checkpoints were approximately 80 MB. These results support the retrospective feasibility of large-scale multimodal analytics with LCAD--RASA, while leaving prospective clinical validation, stronger centre-shift mitigation, and integration of contrastive multimodal pretraining as necessary next steps.
